\section{Exact verification programs}\label{app:certificates}
The three programs below are part of the proof. They use integer or
rational arithmetic for every certification decision. In the C++ programs,
floating-point conversion is used only for diagnostic output.
The mathematical meaning and arithmetic guarantees of the programs are
proved in Sections~\ref{sec:local}--\ref{sec:weighted},
Section~\ref{grid:section}, and
Appendices~\ref{int:appendix} and~\ref{grid:implementation}, and in
the pair derivation below.

\subsection*{Reproduction protocol}
The accompanying directory \texttt{verification/} contains
\texttt{pair.py}, \texttt{interval.cpp}, and \texttt{grid.cpp},
exactly the sources printed below. Run Python with SymPy installed
and with assertions enabled. Compile the C++ sources with a compiler
supporting signed \texttt{\_\_int128} and a plain \texttt{int} type of
at least 32 bits, also with assertions enabled.
For example, from the paper directory:
\begin{lstlisting}[numbers=none]
python verification/pair.py
g++ -std=c++17 -O3 -UNDEBUG verification/interval.cpp -o interval
g++ -std=c++17 -O3 -UNDEBUG verification/grid.cpp -o grid
for j in 0 1 2; do ./interval "$j"; done
for j in -1 $(seq 0 26); do ./grid "$j"; done
\end{lstlisting}
Do not use Python's \texttt{-O} option, and do not define \texttt{NDEBUG}.
Only the displayed argument ranges are part of the protocol.
The bundled runner \nolinkurl{verification/reproduce.py} verifies the source
hashes, runs this suite, checks the expected outputs, and saves the
commands, software versions, exit statuses, and logs in its designated
output directory. Its usage is documented in the accompanying README.

Successful execution of the pair program produces no output.
The interval runs finish with visited-box counts
$414080$, $756096$, and $1004426$, respectively.
The grid run with argument $-1$ stops at iteration $162$.
The other grid stopping iterations are recorded in
Section~\ref{grid:section}.
The certificate sources were replayed with Python~3.12.13,
SymPy~1.14.0, and GCC~13.3.0 on x86\_64 Linux; all 32 mathematical
executions and five environment/compiler checks passed. Separately, the
three interval cases were replayed with \texttt{-O2 -UNDEBUG},
\texttt{-fsanitize=undefined}, and
\texttt{-fno-sanitize-recover=undefined}. All three exited successfully,
with the same progress logs and no sanitizer diagnostics. This supplemental
check is distinct from the ordinary runner protocol and supplies the
execution check on the visit counter discussed in
Appendix~\ref{int:appendix}.

The SHA-256 digests of the three source files, in the above order, are
\begin{quote}\small\ttfamily
\nolinkurl{ab775c3a63fb1d2713bfe058af2dece4ceda99627c17ce18f51f817b69410bba}\\
\nolinkurl{76d1443b1194597ee77749e8030c825f2b9aae46b5c172d9eccf0745078c28e3}\\
\nolinkurl{03c92321efedd0db53b1e91372918dfcd68d459db0223cb160c39df3bd645283}
\end{quote}
These digests identify the exact files; the correctness arguments are
provided in the text, and the successful runs establish the finite
assertions required there.

\subsection{Pair coefficient certificate}\label{cert:pair}
We prove the $a=1/2$ case of Lemma~\ref{loc:pair} for strictly
positive input vectors, using its notation for the input invariants and
pair excesses.  The main-text proof then extends these bounds to every
$0\le a\le1/2$ and to the closed simplex.
Parametrize sorted unnormalized vectors by
\[
 p=(6s+3t+2,3t+2,2),\qquad
 q=(6u+3v+2,3v+2,2),\qquad s,t,u,v\ge0.
\]
This covers every strictly positive sorted vector up to scale: for
$h_0\ge h_1\ge h_2>0$, take
$s=(h_0-h_1)/(3h_2)$ and $t=2(h_1-h_2)/(3h_2)$ to obtain
$p=2h/h_2$.  Put $P=\sum_i p_i$, $Q=\sum_i q_i$, and
$e_i=Q-q_i$, with $E=\sum_i e_i=2Q$.  Then
$e_i/E=(1-q_i/Q)/2$, as required for the coloring edge.
For an unnormalized vector $r$ of total $R$, put
\[
 N_x(r)=\frac{3\sum_i r_i^2-R^2}{2},\qquad
 N_y(r)=\frac{27\prod_i r_i-R^3+3R N_x(r)}2.
\]
The identities for the posterior invariants give
$N_x(r)=R^2x(r/R)$ and $N_y(r)=R^3y(r/R)$.
These are the two numerators returned by \texttt{xy} in the source.

For each of the two choices
$\widetilde e=(e_0,e_1,e_2)$ and $\widetilde e=(e_0,e_2,e_1)$,
let
\[
 r_i^{(j)}=p_i\widetilde e_{i+j},\qquad
 T_j=\sum_i r_i^{(j)},\qquad j=0,1,2,
\]
with indices modulo three.  The likelihood normalizer is
$3T_j/(PE)$; multiplying by the cyclic probability $1/3$ gives weight
$T_j/(PE)$.  Since $\sum_jT_j=PE$, these are probability weights.
The two choices, each followed by the three rotations, enumerate all
six relative permutations.  Repeated representatives at a stabilizer
remain valid uniform group averages.

For each reflection class separately, clear the positive denominator
\begin{equation}\label{loc:denominator}
 D=PE(T_0T_1T_2)^3P^4E^4.
\end{equation}
The three expressions \texttt{out[0]}, \texttt{out[1]}, and
\texttt{out[2]} divided by $D$ are exactly
\[
 \sum_{j=0}^2\frac{N_x(r^{(j)})}{PE T_j},\qquad
 \sum_{j=0}^2\frac{N_y(r^{(j)})}{PE T_j^2},\qquad
 \sum_{j=0}^2\frac{N_x(r^{(j)})^2}{PE T_j^3},
\]
namely the tilted output means of $x,y,x^2$ within that class.
For every call in the source,
\[
 \frac{\texttt{mon}(i,j,k,l)}D=X^i y_0^j z^kB_0^l.
\]
The powers $4-2i-3j$ and $4-2k-3l$ in these calls are nonnegative,
so the cleared objects are polynomials.  Hence \texttt{dx},
\texttt{de}, and \texttt{dn} are $D$ times the class-specific excesses.

The correction \texttt{im}, after division by $D$, is
\begin{equation}\label{loc:reflection}
 I=\frac{27\,\operatorname{odd}(p)\operatorname{odd}(\widetilde e)}
              {64P^3E^3},\qquad
 \operatorname{odd}(p)=4(p_1-p_2)(p_0-p_1)(p_0-p_2).
\end{equation}
Indeed this follows immediately by dividing
$27\operatorname{odd}(p)\operatorname{odd}(\widetilde e)
 [PE(T_0T_1T_2)^3]PE/64$ by~\eqref{loc:denominator}.
Swapping the last two edge coordinates negates $I$.
Although $D$ itself depends on the reflection class, the normalized
correction in~\eqref{loc:reflection} has exactly the indicated sign
change.

To see explicitly what is certified, temporarily write $d_x,d_e,d_n$
for the excesses within one reflection class.  The seven polynomials
checked in that class are the following expressions multiplied by $D$:
\begin{align*}
 &13Xz-2d_n,\\
 &11Xz-d_e,\\
 &250d_e+79Xz,\\
 &-20Xz-40W+40y_0B_0+51S-20(d_x-2I),\\
 &6Xz-6W+y_0B_0+S-(d_e+I),\\
 &2(d_e+I)-12Xz+12W-2y_0B_0+19S,\\
 &-4Xz+44W-14y_0B_0+52S_A+111S_B
       -\{5(d_n+10I)-4(d_e+I)\}.
\end{align*}
The last line, divided by two, bounds
$Cd_n-2d_e+23I$.  All arithmetic in this certificate consists of
integer polynomial addition and convolution, and exact scalar division,
in $\mathbb Z[s,t,u,v]$.  Each of the fourteen polynomials has total
degree $28$ and all its coefficients are nonnegative.  For either
reflection the numbers of nonzero coefficients in the displayed order
are
\[
 13672,\ 13672,\ 13674,\ 13625,\ 13625,\ 13625,\ 13625.
\]
The supplied executable certificate checks every coefficient by the
assertion \texttt{min(f.coeffs())>=0}; absent monomials have coefficient
zero.  Its unchanged-source execution, with Python assertions enabled,
terminates successfully.  Thus these are exact polynomial
nonnegativity statements on the entire nonnegative quadrant, rather
than evaluations on a set of test points.

Divide each class's seven inequalities by its own positive $D$ and
then average the two normalized classes.  All right-hand invariant
terms are unchanged by reflection, and the $I$ terms cancel by
\eqref{loc:reflection}.  This gives precisely
\eqref{loc:coarsepair}--\eqref{loc:kpair} for $a=1/2$.

The complete polynomial calculation is the following source.
\lstinputlisting[language=Python]{../verification/pair.py}

\subsection{Interval positivity certificate}\label{cert:interval}
\lstinputlisting[language=C++]{../verification/interval.cpp}

\subsection{Finite experiment certificate}\label{cert:grid}
\lstinputlisting[language=C++]{../verification/grid.cpp}
